\documentclass[11pt]{article}
\usepackage{amsmath,amssymb,amsthm}
\usepackage[margin=1in]{geometry}
\usepackage{enumitem}

\newcommand{\E}{\mathbb{E}}
\newcommand{\1}{\mathbb{1}}

\title{Math Formulation: GAN Generator with Frozen Permeability
Classifier and GRPO Training --- Continuation}
\author{S Hariharan, J Shrenik Reddy, Ansa N}
\date{}

\begin{document}
\maketitle

\noindent This continues directly from Section 7 of the working draft
(Algorithm Sketch), formalizing the three items listed under
``Open Items to Formalize Next.'' Equation numbers continue from
(8); section numbers continue from 7.

\section{Parameterizing $\pi_\theta(a_t \mid a_{<t}, z)$}
\label{sec:action-param}

Item 1 asked what an ``action'' $a_t$ actually is. This depends on
whether $G_\theta$ is autoregressive (SELFIES) or one-shot
(MolGAN-style atom/bond tensors). Both are formalized below so the
choice can be made explicit in the implementation and the rest of
the math (Eqs.\ 1--8) specialized accordingly.

\subsection{Case A: Autoregressive SELFIES decoding}

Let $\mathcal{V}$ be the SELFIES token vocabulary (including a
terminal \texttt{[EOS]} token) and let $f_\theta$ be a sequence model
(RNN/Transformer) conditioned on the latent $z$. At step $t$,
\begin{equation}
\pi_\theta(a_t \mid a_{<t}, z) \;=\; \operatorname{softmax}\!\big(f_\theta(a_{<t}, z)\big)_{a_t},
\qquad a_t \in \mathcal{V},
\label{eq:selfies-step}
\end{equation}
and decoding terminates at the first $t = T$ with $a_T = \texttt{[EOS]}$
(or a fixed maximum length $T_{\max}$). $T$ is therefore a random
variable determined by the rollout itself, not a hyperparameter.
The molecule is obtained by the deterministic SELFIES$\to$graph
parser, $m = \mathrm{Parse}(a_1,\dots,a_T)$, so $\pi_\theta(m)$ in
Eq.~1 is exactly the product over the \emph{realized} $T$:
\begin{equation}
\pi_\theta(m) = \prod_{t=1}^{T} \pi_\theta(a_t \mid a_{<t}, z), \qquad
\log \pi_\theta(m) = \sum_{t=1}^{T} \log \pi_\theta(a_t \mid a_{<t}, z).
\label{eq:selfies-logprob}
\end{equation}
Because $\mathrm{Parse}$ is guaranteed syntactically valid by
construction of the SELFIES grammar (every token string maps to
\emph{some} molecule), $\mathrm{Valid}(m)$ in Section~4 is checking
chemical/RDKit sanitization validity, not SELFIES syntax validity ---
these are different failure modes and should not be conflated in the
implementation.

The per-token ratio $r_t(\theta)$ from Eq.~2 is unchanged and applies
token-by-token; the clipped objective (Eq.~7) is naturally a sum
over $t = 1,\dots,T$ per sampled molecule.

\subsection{Case B: One-shot atom/bond tensor (MolGAN-style)}

Here $G_\theta$ maps $z$ directly to two tensors in a single forward
pass: a node/atom-type tensor $X \in \mathbb{R}^{N_{\max}\times
|\mathcal{A}|}$ (logits over atom types per node slot, including a
``no atom'' padding type) and an edge/bond-type tensor
$E \in \mathbb{R}^{N_{\max}\times N_{\max}\times|\mathcal{B}|}$
(logits over bond types per node pair, including ``no bond''), where
$N_{\max}$ is a fixed maximum atom count. There is no notion of a
sequential decision $a_1,\dots,a_T$; instead there is a single joint
action
\begin{equation}
a = (x_1,\dots,x_{N_{\max}},\, e_{12}, e_{13},\dots, e_{(N_{\max}-1)N_{\max}}),
\end{equation}
with each $x_i \in \mathcal{A}$ sampled from
$\operatorname{Categorical}(\operatorname{softmax}(X_i))$ and each
$e_{ij} \in \mathcal{B}$ sampled from
$\operatorname{Categorical}(\operatorname{softmax}(E_{ij}))$,
independently conditioned on $z$ (the joint distribution over the
tensor entries factorizes given $z$, even though the resulting graph
is not independent of itself once parsed and sanitized). The action
probability is then
\begin{equation}
\pi_\theta(a \mid z) = \prod_{i=1}^{N_{\max}} \pi_\theta(x_i \mid z)
\;\cdot\! \prod_{1\le i<j\le N_{\max}} \pi_\theta(e_{ij}\mid z).
\label{eq:molgan-joint}
\end{equation}
In the notation of Eq.~1, this is the degenerate $T=1$ case: the
``sequence'' collapses to one joint categorical draw, and $\pi_\theta(m)
= \pi_\theta(a\mid z)$ directly (no autoregressive conditioning on
previous atoms/bonds). Consequently the ratio $r_t(\theta)$ in
Eq.~2 is not indexed by $t$ at all; it becomes a single per-molecule
ratio
\begin{equation}
r(\theta) = \frac{\pi_\theta(a\mid z)}{\pi_{\theta_{\mathrm{old}}}(a\mid z)},
\label{eq:molgan-ratio}
\end{equation}
and the clipped surrogate (Eq.~7) drops its inner $\E_t[\cdot]$,
reducing to $L^{\mathrm{CLIP}}_i(\theta) = \min\big(r(\theta) A_i,\,
\mathrm{clip}(r(\theta), 1\pm\varepsilon)\, A_i\big)$ evaluated once
per sampled molecule $i$, not once per token.

\paragraph{Note on differentiability.} MolGAN-style generators are
often trained with a Gumbel-softmax relaxation so that samples remain
differentiable and the discriminator loss can backpropagate through
$X, E$ directly. That reparameterization path is \emph{not} needed
here: because discrete molecule construction already routes through
the non-differentiable RDKit sanitization step (Section 2), the whole
pipeline is trained with the REINFORCE/PPO score-function estimator
via $R(m)$, not via a straight-through gradient. If a Gumbel-softmax
relaxation is used at all, it should be restricted to stabilizing
$D_\phi$'s own adversarial loss in Eq.~3--4 (which is a separate,
directly differentiable objective), and kept out of the $\theta$
update path to avoid mixing two different gradient estimators for
the same parameters.

\subsection{Which case governs $r_t(\theta)$ downstream}

Eqs.~7--8 should be read as written (indexed by $t$) under Case A,
and with $t$ collapsed to a single joint action under Case B via
Eq.~\eqref{eq:molgan-ratio}. This is the concrete fork that Section
7's Algorithm Sketch (step 2, ``Generate molecules'') needs to commit
to before implementation; nothing else in Sections 3--6 depends on
the choice.

\section{Fixed vs.\ Annealed Reward Weights}
\label{sec:weight-schedule}

Item 2 asked whether $w_D, w_C, w_Q, w_S$ in Eq.~5 are fixed or
annealed. Write $k$ for the generator training step (the outer loop
counter in Section 7, step 9).

\subsection{Fixed-weight baseline}
The simplest choice keeps $w_D, w_C, w_Q, w_S$ constant for all $k$,
subject only to $w_D, w_C, w_Q, w_S \ge 0$. This is the version
implicitly used in Eq.~5 and is the right default to implement first,
since it has no extra hyperparameters beyond the weights themselves.

\subsection{Annealed schedule}
Motivated by the concern in the draft (establish validity/realism
first, then shift weight to $C$), define each weight as a function of
$k$ interpolating between an initial and final value over an
annealing horizon $k_{\mathrm{anneal}}$:
\begin{equation}
w_\bullet(k) = w_\bullet^{(0)} + \big(w_\bullet^{(f)} - w_\bullet^{(0)}\big)\cdot
\lambda(k), \qquad
\lambda(k) = \min\!\Big(\tfrac{k}{k_{\mathrm{anneal}}},\, 1\Big),
\label{eq:linear-anneal}
\end{equation}
for $\bullet \in \{D, C, Q, S\}$, with e.g.\ $w_D^{(0)} \gg w_C^{(0)}$
and $w_D^{(f)} \ll w_C^{(f)}$ to realize the ``validity first, then
permeability'' curriculum. $\lambda(k)$ can equally be a cosine
schedule, $\lambda(k) = \tfrac12\big(1-\cos(\pi \min(k/k_{\mathrm{anneal}},1))\big)$,
if smoother transitions are preferred; the choice does not change any
other equation, since $w_\bullet(k)$ only enters Eq.~5 as a
per-step constant.

\subsection{Validity-gated curriculum switch}
An alternative to a fixed $k_{\mathrm{anneal}}$ is to gate the
transition on measured performance rather than wall-clock steps.
Let $\rho_{\mathrm{valid}}(k) = \frac{1}{N}\sum_{i=1}^N
\1[\mathrm{Valid}(m_i)]$ be the fraction of valid molecules in step
$k$'s sampled group. Replace $\lambda(k)$ in
Eq.~\eqref{eq:linear-anneal} with
\begin{equation}
\lambda(k) = \1\!\big[\bar\rho_{\mathrm{valid}}(k) \ge \tau\big],
\qquad
\bar\rho_{\mathrm{valid}}(k) = \frac{1}{W}\sum_{j=k-W+1}^{k} \rho_{\mathrm{valid}}(j),
\end{equation}
a hard switch once the $W$-step moving average of validity rate
exceeds threshold $\tau$ (e.g.\ $\tau = 0.9$). This ties the
curriculum to the generator's actual competence at producing
sanitizable molecules rather than to an arbitrarily chosen step
count, at the cost of one more hyperparameter pair $(\tau, W)$. Either
schedule is a drop-in replacement for the constant weights in Eq.~5;
no other equation in Sections 4--8 needs to change.

\section{Discriminator/Generator Update Ratio $n_D$ and Reward Non-Stationarity}
\label{sec:nD-nonstationarity}

Item 3 asked for a formal treatment of $n_D$ (Section 7, step 8) and
its effect on the reward's non-stationarity, since $D_\phi$ --- unlike
the frozen classifier $C$ --- is itself being trained.

\subsection{Why this is a genuine non-stationarity, and why $C$ is exempt}
Write $R_k(m)$ for the reward function actually used at generator
step $k$, i.e.\ Eq.~5 with $D_\phi$ replaced by $D_{\phi_k}$, the
discriminator parameters after its $k$-th block of $n_D$ updates.
Because $\phi_k$ changes with $k$, $R_k$ is a genuinely time-varying
reward: the same molecule $m$ can receive a different $w_D
D_{\phi_k}(m)$ term at step $k$ than it would have at step $k' \ne k$,
even though nothing about $m$ changed. This is exactly the standard
GAN dynamics problem transplanted into the reward function. By
contrast $C$ is frozen ($\nabla_\phi C = 0$, Section 4), so the
$w_C C(m)$ term is stationary in $k$ by construction --- it is only
the $D_\phi$ term that introduces non-stationarity.

\subsection{Two-timescale formalization}
The standard way to make convergence arguments tractable here
(following the two-timescale update rule (TTUR) analysis used for
GANs) is to treat $\phi$ as evolving on a \emph{faster} timescale than
$\theta$: choose $n_D \ge 1$ discriminator gradient steps per
generator step (or a discriminator learning rate $\eta_\phi \gg
\eta_\theta$) such that, on the timescale over which $\theta$ changes
appreciably, $\phi_k$ can be treated as approximately equilibrated to
the current generator distribution $p_{\theta_k}$:
\begin{equation}
\phi_k \;\approx\; \phi^\star(\theta_k) := \arg\min_\phi \; L_D(\phi;\, \theta_k),
\label{eq:ttur}
\end{equation}
with $L_D$ as in Eq.~3 (or Eq.~4). Under this quasi-static assumption,
$R_k(m) \approx R(m;\theta_k)$ becomes a (slowly) $\theta$-dependent
reward rather than an arbitrarily-moving target, which is the
condition under which the GRPO update in Section 5--6 can be
interpreted as ordinary policy-gradient ascent on a well-defined,
slowly-drifting objective at each step, rather than pursuing a
reward that has no stationary point. This is an assumption, not a
guarantee: it degrades as $n_D \to 1$ (discriminator barely updated
before the next generator step) and improves as $n_D$ grows, at the
usual cost of more discriminator compute per generator step.

\subsection{Effect on the group-relative advantage}
The advantage $A_i$ (Eq.~6) is computed \emph{within} a single step
$k$, using $\mu_{R_k}, \sigma_{R_k}$ over the current group only ---
so intra-group ranking of candidates is unaffected by $D_\phi$ drift
across steps: every candidate in the group is scored by the same
$D_{\phi_k}$ snapshot. The non-stationarity instead shows up
\emph{across} steps, in two ways worth tracking empirically rather
than in the group statistics themselves:
\begin{enumerate}[nosep]
  \item \textbf{Signal magnitude drift.} As $D_\phi$ sharpens over
  training (its output distribution typically becomes less entropic,
  pushing $D_\phi(m)$ toward $\{0,1\}$), the \emph{variance} of the
  $w_D D_\phi(m)$ term across a group can grow or shrink relative to
  the fixed-variance $w_C C(m)$ term, changing the effective relative
  weight the two terms have on $A_i$ even if $w_D, w_C$ are held
  literally constant (Section 9.1). This is worth logging
  ($\mathrm{Var}_i[D_{\phi_k}(m_i)]$ per step) as a diagnostic.
  \item \textbf{Off-policy staleness within an update.} Section 6
  reuses the same sampled group for several PPO epochs/minibatches
  before resampling (standard PPO reuse pattern). If $\phi$ is also
  updated during that same window (i.e.\ $n_D$ discriminator steps
  interleaved with the $K$ generator epochs of step 7 in Section 7),
  the reward used to compute $A_i$ at the start of the window becomes
  stale relative to $\phi$ by the end of it. The safest choice is to
  freeze $\phi$ for the duration of a generator update window (i.e.
  perform all $n_D$ discriminator steps either strictly before or
  strictly after the $K$-epoch PPO block, never interleaved with it),
  so that $A_i$ is computed once against a single fixed $\phi$
  snapshot per group, consistent with the frozen-$C$ treatment in
  Section 4.
\end{enumerate}

\subsection{Practical implication for choosing $n_D$}
Combining 10.2 and 10.3: $n_D$ should be chosen large enough that
Eq.~\eqref{eq:ttur} is a reasonable approximation (discriminator
meaningfully ahead of the generator, as in standard WGAN-GP
practice, e.g.\ $n_D \in [1,5]$), but the discriminator update block
should be scheduled \emph{outside} the PPO reuse window from Section
6 rather than concurrently with it, so that $R_k$ --- and hence
$A_i$ --- is well-defined as a single fixed function for the entire
duration over which it is used to compute the clipped objective in
Eq.~7.

\end{document}
